\pagebreak
\section{Описание}

\subsection{Корпус}
Токенизатор проверяется на корпусе первой лабораторной работы --- 30\,000 случайных
статей русской Википедии (поле \texttt{text} в JSON-файлах, очищенный текст без
таблиц и сносок). Объём текста --- 231{,}2 МБ. Википедия хорошо подходит для
проверки: в ней много сокращений (\enquote{г.}, \enquote{н.\,э.}, \enquote{т.\,д.}),
инициалов (\enquote{А.\,С. Пушкин}), чисел всех видов (даты, десятичные дроби,
числа с пробелами-разделителями), составных слов (\enquote{Ростов-на-Дону},
\enquote{северо-запад}), а также встречаются латиница, ссылки и адреса почты.

\subsection{Общее устройство}
Токенизатор --- модуль \texttt{main.py} на Python 3 без внешних зависимостей
(только стандартная библиотека). Вся логика --- один регулярный
шаблон с именованными группами и функция \texttt{tokenize(text)}, которая
возвращает список пар \texttt{(вид, токен)}. Шаблон просматривает текст слева
направо, на каждой позиции пробует варианты \emph{в порядке приоритета} и берёт
первый подошедший; поэтому более специфичные правила (ссылка, инициалы,
сокращение, число) стоят раньше общих (слово). Всё, что не подошло ни к одному
правилу (пробелы, подчёркивание), пропускается.

\subsection{Правила токенизации}
\begin{longtable}{|p{2.4cm}|p{6.2cm}|p{5.8cm}|}
\hline
\textbf{Вид} & \textbf{Правило} & \textbf{Примеры} \\ \hline
\texttt{url}, \texttt{email} & \texttt{http(s)://} или \texttt{www.} до пробела; знаки
препинания в конце не входят в ссылку; адрес почты & \texttt{www.eol.org},
\texttt{https://doi.org/10.1/x} \\ \hline
\texttt{initials} & заглавная буква с точкой (одна или две подряд) перед
фамилией; римские цифры \texttt{I}, \texttt{V}, \texttt{X} исключены, иначе
\enquote{Пётр I. Он} склеивается & \enquote{А.\,С. Пушкин} $\to$ \texttt{а.с.}, \texttt{пушкин} \\ \hline
\texttt{abbr} & \enquote{т.\,д.}, \enquote{т.\,е.}, \enquote{н.\,э.}; цепочки
\enquote{буква-точка}; слова из словаря с точкой (\enquote{г.}, \enquote{гг.},
\enquote{руб.}, \enquote{англ.}, \dots); \enquote{в.}/\enquote{г.} после числа или римской
цифры & \texttt{т.е.}, \texttt{н.э.}, \texttt{г.}, \texttt{руб.}, \texttt{xx} + \texttt{в.} \\ \hline
\texttt{number} & целые, десятичные (\texttt{3,14}), даты (\texttt{12.03.2020}),
время (\texttt{10:30}); группы по три цифры через пробел (\texttt{10 000}) & \texttt{3,14}, \texttt{10 000} \\ \hline
\texttt{compound} & буквы и цифры через дефис; диапазон из одних цифр
(\texttt{1990-2000}) разбивается на два числа & \texttt{ростов-на-дону}, \texttt{кто-то}, \texttt{f-16}, \texttt{1990-е} \\ \hline
\texttt{word} & последовательность букв и цифр, апостроф между буквами
не разрывает слово & \texttt{пушкин}, \texttt{o'neil}, \texttt{3d} \\ \hline
\texttt{punct} & любой другой знак; по умолчанию отбрасывается, возвращается
при \texttt{punct=True} & \texttt{,} \texttt{.} \texttt{(} \\ \hline
\end{longtable}

\subsection{Нормализация}
\begin{itemize}
    \item Приведение к нижнему регистру (\texttt{lower=True}), замена \enquote{ё} на
    \enquote{е} (\texttt{yo=True}): в Википедии одно и то же слово пишется обоими
    способами, и без замены оно распадается на два термина.
    \item Удаляются мягкий перенос (U+00AD) и знак ударения (U+0301):
    слово с ударением на второй букве (\enquote{Глау}) иначе распадается на два токена.
    \item Пробелы внутри сокращений убираются: \enquote{т. е.} $\to$ \texttt{т.е.}
    \item Составные слова по умолчанию остаются одним токеном; параметр
    \texttt{split\_}\allowbreak\texttt{compounds=True} разбивает их на части (\enquote{ростов},
    \enquote{на}, \enquote{дону}) --- это нужно, если поиск по части слова важнее
    точного совпадения.
\end{itemize}

\subsection{Запуск}
\begin{lstlisting}[language=bash]
python main.py                        # whole corpus, statistics
python main.py --limit 1000           # first 1000 documents
python main.py --text "text to tokenize"
python test_main.py                   # unit checks
\end{lstlisting}

\subsection{Пример работы}
Предложение: \enquote{В 1823 г. А.\,С. Пушкин жил в Ростове-на-Дону, т.\,е. на юге:
10 000 руб., 3,14\% на 12.03.2020.} Сравнение с наивным разбиением по
\texttt{\textbackslash w+} (всё, что не буква и не цифра, --- граница):

\begin{longtable}{|p{3.2cm}|p{11.1cm}|}
\hline
\textbf{Метод} & \textbf{Токены} \\ \hline
\texttt{\textbackslash w+} (24) & в, 1823, г, а, с, пушкин, жил, в, ростове, на, дону, т, е, на, юге,
10, 000, руб, 3, 14, на, 12, 03, 2020 \\ \hline
Токенизатор (16) & в, 1823, г., а.с., пушкин, жил, в, ростове-на-дону, т.е., на, юге,
10 000, руб., 3,14, на, 12.03.2020 \\ \hline
\end{longtable}

Наивный метод рвёт \enquote{а.\,с.}, \enquote{т.\,е.} на одиночные буквы (шум: токены
\enquote{а}, \enquote{с}, \enquote{т}, \enquote{е} совпадают с предлогами и союзами),
режет \enquote{10 000} и дату на несколько чисел, а \enquote{3,14} превращает в
два числа. Токенизатор сохраняет каждую из этих единиц как один токен.

\subsection{Результаты на корпусе}
Прогон по всем 30\,000 документам (один поток, Python 3.13):

\begin{longtable}{|p{8.2cm}|p{6.1cm}|}
\hline
\textbf{Показатель} & \textbf{Значение} \\ \hline
Документов / объём текста & 30\,000 / 231{,}2 МБ \\ \hline
Токенов (наш токенизатор) & 17\,765\,479 \\ \hline
Токенов (\texttt{\textbackslash w+}) & 18\,287\,320 \\ \hline
Токенов (с \texttt{split\_}\allowbreak\texttt{compounds}) & 18\,072\,387 \\ \hline
Уникальных токенов (наш / \texttt{\textbackslash w+} / с разбиением) & 978\,389 / 896\,072 / 912\,413 \\ \hline
Токенов, встречающихся один раз (наш) & 519\,230 (53\%) \\ \hline
Средняя длина токена & 6{,}08 символа \\ \hline
Время токенизации & 35{,}5 с \\ \hline
Скорость & 6{,}2 МБ/с (6368 КБ/с), 500\,тыс.\ токенов/с \\ \hline
\end{longtable}

Распределение по видам токенов:

\begin{longtable}{|l|r|r|}
\hline
\textbf{Вид} & \textbf{Токенов} & \textbf{Доля} \\ \hline
\texttt{word} & 16\,006\,198 & 90{,}2\% \\ \hline
\texttt{number} & 1\,231\,352 & 6{,}9\% \\ \hline
\texttt{compound} & 291\,642 & 1{,}6\% \\ \hline
\texttt{abbr} & 112\,483 & 0{,}6\% \\ \hline
\texttt{initials} & 95\,304 & 0{,}5\% \\ \hline
\texttt{url} & 2\,429 & $<$0{,}1\% \\ \hline
\texttt{email} & 9 & $<$0{,}1\% \\ \hline
\end{longtable}

Распределение по длине токена: 1 символ --- 11{,}0\%, 2--3 --- 14{,}5\%,
4--7 --- 42{,}4\%, 8 и больше --- 32{,}1\%.

Самые частые токены разных видов:
\begin{itemize}
    \item составные: \enquote{из-за}, \enquote{1-й}, \enquote{2-й}, \enquote{санкт-петербург},
    \enquote{генерал-майор}, \enquote{нью-йорк};
    \item сокращения: \enquote{англ.}, \enquote{г.}, \enquote{см.}, \enquote{гг.},
    \enquote{ред.}, \enquote{н.э.};
    \item инициалы: \enquote{а.а.}, \enquote{в.а.}, \enquote{в.в.}, \enquote{в.и.};
    \item ссылки: \enquote{www.youtube.com}, \enquote{www.sports-reference.com}.
\end{itemize}

\subsection{Тестирование}
Файл \texttt{test\_main.py} содержит 18 проверок правил (пунктуация, \enquote{ё},
ударение, составные слова, тире с пробелами, диапазоны, числа с разделителями,
сокращения, инициалы, римские цифры, ссылки и почта) и проверки параметров
\texttt{split\_}\allowbreak\texttt{compounds}, \texttt{lower}, \texttt{punct}. Все проверки
проходят (\texttt{python test\_main.py} печатает \texttt{ok}).
